\begin{examplebox}
\textbf{例1：代数直接求比值。} 椭圆 $x^2/25+y^2/16=1$ 中
$a=5,b=4,c=3$，故 $\boxed{e=3/5}$。
回验：右顶点至右焦点距离为 $2$，至对应准线 $x=25/3$ 距离为 $10/3$，比为 $3/5$。

\medskip
\textbf{例2：已知离心率反求双曲线。}
中心在原点、焦点在 $x$ 轴，实轴长为 $8$，$e=5/4$。
先取 $a=4$，再由 $c=ae=5$ 得 $b^2=25-16=9$，所以
\[
 \boxed{\frac{x^2}{16}-\frac{y^2}{9}=1}.
\]
回验：$c/a=5/4>1$，实轴长 $2a=8$，焦点确在 $x$ 轴。

\medskip
\textbf{例3：几何比值转为方程，并核查点存在。}
中心在原点、焦点在 $x$ 轴的椭圆长轴长 $20$；某点 $P$ 满足到右焦点距离 $4$、到对应准线距离 $5$。
由定义 $e=4/5$，$a=10,c=8$，故 $b^2=100-64=36$，答案为
\[
 \boxed{\frac{x^2}{100}+\frac{y^2}{36}=1}.
\]
\begin{center}\begin{tikzpicture}[font=\small,>=stealth]
\draw[black!30,rounded corners=2pt] (0,0) rectangle (14.4,8.2);
\node[anchor=west,font=\small\bfseries] at (.5,7.72) {例题核查：同一个点，配错准线就不再得到原离心率};

\begin{scope}[shift={(7.2,3.8)},scale=.38]
\draw[blue,very thick] (0,0) ellipse (10 and 6);
\draw[black!35] (-13,0)--(13,0);
\draw[violet,thick] (-12.5,-6.5)--(-12.5,6.5);
\draw[violet,thick] (12.5,-6.5)--(12.5,6.5);
\draw[orange!85!black,very thick] (8,0)--(7.5,3.9686269666);
\draw[black!60,very thick] (7.5,3.9686269666)--(12.5,3.9686269666);
\draw[violet,dashed] (-12.5,3.9686269666)--(7.5,3.9686269666);
\fill (8,0) circle (3pt) node[below left=6pt] {$F_R(8,0)$};
\fill (7.5,3.9686269666) circle (3pt) node[above=12pt] {$P$};
\node[left=7pt] at (7.8,1.8) {$PF_R=4$};
\node[above=5pt] at (10,3.9686269666) {$PH_R=5$};
\node at (-4,6.8) {$PH_L=20$};
\end{scope}
\node at (2.45,6.85) {$l_L:x=-25/2$};
\node at (11.95,6.85) {$l_R:x=25/2$};
\node at (7.2,1) {正确配对：$4/5$；错误配对：$4/20=1/5$};

\node[anchor=west] at (.7,.55) {$P=(15/2,3\sqrt7/2)$ 在 $x^2/100+y^2/36=1$ 上；焦点与准线须同侧对应};
\end{tikzpicture}\end{center}

回验：右准线为 $x=25/2$，$d=5$ 给出 $x_P=15/2$；
代回方程得 $y_P=\pm3\sqrt7/2$，确有两点。
再算 $PF_R^2=(-1/2)^2+(3\sqrt7/2)^2=16$，故 $PF_R=4$。
换成左准线会得到 $4/20=1/5$，不等于本椭圆离心率。
\end{examplebox}
\clearpage

\begin{examplebox}
\textbf{例4：不背准线，也能推离心率（椭圆两顶点法）。}
某横向椭圆的半长轴 $a=5$、半焦距 $c=3$，已知该椭圆满足焦准距离比恒定。
令对应的右准线为 $x=t$（$t>5$）。
两顶点对\textbf{同一个右焦点} $(3,0)$ 的距离分别为 $2$、$8$；
到右准线的距离分别为 $t-5$、$t+5$。于是
\[
 \frac2{t-5}=\frac8{t+5}=e
 \quad\Longrightarrow\quad
 2(t+5)=8(t-5)
 \quad\Longrightarrow\quad t=\frac{25}{3}.
\]
故 $e=2/(25/3-5)=\boxed{3/5}$。
\textbf{更快：} 不必先解 $t$，把 $t-5=2/e$、$t+5=8/e$ 相减，
直接得 $10=6/e$，即 $e=3/5$。

\medskip
\textbf{例5：同一招应对双曲线的符号。}
双曲线 $x^2/9-y^2/16=1$，$a=3,b=4,c=5$。
设右准线 $x=t$（$0<t<3$）。右、左顶点到右焦点的距离为 $2,8$，
到右准线距离为 $3-t,3+t$，所以
\[
 \frac{2}{3-t}=\frac{8}{3+t}=e
 \quad\Longrightarrow\quad
 \begin{cases}3-t=2/e,\\ 3+t=8/e\end{cases}
 \quad\Longrightarrow\quad
 6=10/e.
\]
即 $\boxed{e=5/3}$；回代得 $t=9/5$。
\textbf{自检：} 椭圆准线在顶点外，双曲线准线在两顶点内；
但这两个例子都始终使用\textbf{右焦点配右准线}。
\end{examplebox}
\clearpage
